% Préambule commun aux manuels de Tutti (français et anglais).
% Compilation : pdflatex manuel-fr.tex (deux fois), idem manuel-en.tex, depuis docs/src/.
\documentclass[11pt,a4paper]{article}
\usepackage[utf8]{inputenc}
\usepackage[T1]{fontenc}
\usepackage{lmodern}
\usepackage{microtype}
\usepackage[a4paper,margin=2cm,headheight=24pt]{geometry}
\usepackage{graphicx}
\usepackage[table]{xcolor}
\usepackage[most]{tcolorbox}
\usepackage{enumitem}
\usepackage{booktabs}
\usepackage{array}
\usepackage{tabularx}
\usepackage[nobottomtitles*]{titlesec}
\usepackage{needspace}
\usepackage{fancyhdr}
\usepackage{fontawesome5}
\usepackage{etoolbox}
\DeclareUnicodeCharacter{2192}{\ensuremath{\rightarrow}}
\DeclareUnicodeCharacter{2191}{\ensuremath{\uparrow}}
\DeclareUnicodeCharacter{2193}{\ensuremath{\downarrow}}
\DeclareUnicodeCharacter{25B2}{\ensuremath{\blacktriangle}}
\DeclareUnicodeCharacter{25BC}{\ensuremath{\blacktriangledown}}
\usepackage{amssymb}

% Palette de l'application : ivoire, bordeaux, or
\definecolor{bordeaux}{HTML}{7A1F2B}
\definecolor{bordeauxl}{HTML}{F6E7E9}
\definecolor{or}{HTML}{B8860B}
\definecolor{orl}{HTML}{FBF3DC}
\definecolor{sapin}{HTML}{2C7A55}
\definecolor{sapinl}{HTML}{E2F1E9}
\definecolor{sky}{HTML}{1C6DB4}
\definecolor{skyl}{HTML}{E6F0F8}
\definecolor{summit}{HTML}{6B7A8F}
\definecolor{night}{HTML}{2A2420}
\definecolor{ink}{HTML}{1F2430}
\AtBeginDocument{\color{ink}}

\titleformat{\section}
  {\Large\bfseries\sffamily\color{bordeaux}}
  {\makebox[1.6em][l]{\thesection}}{0.3em}{}
  [{\vspace{1pt}\color{bordeaux!35}\titlerule[1.3pt]}]
\titlespacing*{\section}{0pt}{14pt}{7pt}
\titleformat{\subsection}{\large\bfseries\sffamily\color{night}}{\thesubsection}{0.5em}{}
\pretocmd{\section}{\needspace{10\baselineskip}}{}{}
\pretocmd{\subsection}{\needspace{7\baselineskip}}{}{}
\widowpenalty=10000
\clubpenalty=10000
\raggedbottom
\setlength{\parindent}{0pt}
\setlength{\parskip}{6pt}
\setlist[enumerate]{itemsep=3pt,leftmargin=6mm}
\setlist[itemize]{itemsep=2pt,leftmargin=5mm}

\pagestyle{fancy}
\fancyhf{}
\renewcommand{\headrulewidth}{0.4pt}
\renewcommand{\headrule}{\hbox to\headwidth{\color{bordeaux!40}\leaders\hrule height \headrulewidth\hfill}}
\fancyfoot[C]{\sffamily\small\color{summit}\thepage}
\renewcommand{\sectionmark}[1]{\markboth{#1}{}}

\newcommand{\menu}[1]{\textsf{\textbf{\textcolor{bordeaux}{#1}}}}
\newcommand{\touche}[1]{\tcbox[on line,boxsep=0pt,left=2pt,right=2pt,top=1.5pt,bottom=1.5pt,
  colback=white,colframe=summit,boxrule=0.5pt,arc=1.5pt]{\footnotesize\textsf{#1}}}
\newcommand{\fichier}[1]{\texttt{\small #1}}

\newcommand{\boite}[5]{\newtcolorbox{#1}{enhanced,unbreakable,colback=#2,colframe=#3,
  coltitle=white,fonttitle=\bfseries\sffamily\small,boxrule=0.9pt,arc=2.5pt,
  left=2.5mm,right=2.5mm,top=1.2mm,bottom=1.2mm,title={#4~~#5}}}
\newcommand{\capture}[2][1]{\par\needspace{8\baselineskip}\begin{center}\setlength{\fboxsep}{0pt}\setlength{\fboxrule}{0.6pt}%
  \fcolorbox{summit!50}{white}{\includegraphics[width=\dimexpr#1\linewidth-1.3pt\relax]{#2}}\end{center}}

\usepackage[english]{babel}
\usepackage[colorlinks=true,linkcolor=bordeaux,urlcolor=bordeaux]{hyperref}
\hypersetup{pdftitle={Tutti — user guide},pdfauthor={},pdfcreator={Tutti}}
\boite{astucebox}{orl}{or!85!black}{\faLightbulb}{Good to know}
\boite{securitebox}{sapinl}{sapin}{\faShield*}{Your files are safe}
\newcommand{\astuce}[1]{\par\needspace{5\baselineskip}\begin{astucebox}#1\end{astucebox}}
\fancyhead[L]{\sffamily\small\color{bordeaux}\textbf{TUTTI — USER GUIDE}}
\fancyhead[R]{\sffamily\small\color{summit}\leftmark}
\input{version}

\begin{document}

\begin{titlepage}
\centering
\vspace*{2.2cm}
\includegraphics[width=3.6cm]{../../logo_1024.png}\par
\vspace{8mm}
{\sffamily\Huge\bfseries\color{bordeaux} TUTTI\par}
\vspace{3mm}
{\sffamily\LARGE\color{ink} Merge, convert, touch up and split PDF files\par}
\vspace{2mm}
{\sffamily\large\color{summit} User guide \textbullet{} version \versiontutti\par}
\vspace{12mm}
\begin{tcolorbox}[enhanced,width=0.8\linewidth,colback=orl,colframe=or,arc=4pt,boxrule=1.2pt,left=5mm,right=5mm,top=3mm,bottom=3mm]
\sffamily\centering
\textit{Tutti}: “all together”, in music.\\[2mm]
\raggedright
Four tabs, in the order of a job:\\[1mm]
\textbf{Merge} \quad → \quad \textbf{Pages} \quad → \quad \textbf{Finishing} \qquad and \quad \textbf{Split}
\end{tcolorbox}
\vfill
{\footnotesize\color{summit}\sffamily Windows \textbullet{} macOS \textbullet{} Linux\par}
\end{titlepage}

\tableofcontents
\newpage

\section{What Tutti does}

Tutti replaces websites such as “iLovePDF”: everything happens on your computer, nothing is
sent over the Internet. Each tab of the window does one job:

\begin{center}
\renewcommand{\arraystretch}{1.4}
\begin{tabularx}{\linewidth}{>{\raggedright\arraybackslash}X >{\sffamily\bfseries\color{bordeaux}}l c}
\toprule
\textbf{I want to…} & \textbf{\color{ink}Tab} & \textbf{Page} \\
\midrule
merge several PDF files into one (full score + parts, one bookmark per part) & Merge and convert & \pageref{sec:assembler} \\
turn a photo, a scan, a Word, Excel or PowerPoint file into a PDF & Merge and convert & \pageref{sec:convertir} \\
make a booklet to fold and staple (A5 on A4, or A4 on A3) & Merge and convert & \pageref{sec:options} \\
straighten a page scanned crooked or sideways & Pages & \pageref{sec:pages} \\
remove a blank or duplicate page, change the page order & Pages & \pageref{sec:pages} \\
number the pages, write a title at the top of every page & Finishing & \pageref{sec:finitions} \\
make a PDF small enough to send by e-mail & Finishing & \pageref{sec:finitions} \\
make one file per page, or extract pages & Split & \pageref{sec:decouper} \\
\bottomrule
\end{tabularx}
\end{center}

\begin{securitebox}
\begin{itemize}
\item The original files are \textbf{never modified}: Tutti reads them and only writes the new PDF.
\item Pages are copied \textbf{as they are}, with no loss of quality and no “re-printing”.
\item Tutti refuses to write over one of the files being merged.
\item A power cut while writing never leaves a half-written PDF behind.
\end{itemize}
\end{securitebox}

\section{Installing Tutti}

Download the latest version from the project's \textbf{Releases} page, picking the file for your
system:

\begin{center}
\renewcommand{\arraystretch}{1.3}
\begin{tabular}{>{\bfseries}l l}
\toprule
Windows & \fichier{Tutti-Setup-X.Y.Z.exe} \\
macOS (Apple M1, M2… chips) & \fichier{Tutti-X.Y.Z-macos-arm64.dmg} \\
macOS (Intel processor) & \fichier{Tutti-X.Y.Z-macos-x86\_64.dmg} \\
Linux & \fichier{Tutti-X.Y.Z-linux-x86\_64.tar.gz} (or \fichier{-arm64}) \\
\bottomrule
\end{tabular}
\end{center}

\textbf{No administrator password is needed}: Tutti installs into your own user profile.

\subsection*{Windows}
Run \fichier{Tutti-Setup-X.Y.Z.exe}. On the very first launch Windows may show “Windows protected
your PC”, because the installer is not signed: click \menu{More info}, then \menu{Run anyway}.
This warning appears only once.

\subsection*{macOS}
Open the \fichier{.dmg} file and drag Tutti into \menu{Applications}. On the first launch macOS
refuses to open an unsigned application: \textbf{right-click} Tutti, choose \menu{Open}, and confirm.

\subsection*{Linux}
Unpack the archive and run \fichier{installer.sh}. Tutti installs into
\fichier{\textasciitilde/.local/share} and appears in the applications menu.

\section{First launch}

Tutti opens in your system's language, French or English. The \menu{Language} menu lets you
force either one.

You are asked once whether to \textbf{send a report automatically if Tutti runs into a problem}.
The report contains the program and system versions and the error; your file names and user name
are hidden. You can change your mind at any time in the \menu{Help} menu.

\capture[0.82]{../capture_vide-en.png}

\section{Merging PDF files, step by step}\label{sec:assembler}

\subsection{Adding the files}
\begin{itemize}
\item \textbf{Drag} the PDF files from Explorer (or the Finder) into the large list;
\item or drag \textbf{a whole folder}: every PDF it contains is added, sorted by name;
\item or use \menu{Add files…} \ (\touche{Ctrl}+\touche{O}) and \menu{Add a folder…}.
\end{itemize}
Each file shows its page count and size. A file already in the list is not added twice.

\subsection{Putting them in order}
The order of the list is the order of the merged PDF.
\begin{itemize}
\item \textbf{Drag} a file up or down inside the list;
\item or select it and use \menu{▲ Move up} / \menu{▼ Move down}
      (\touche{Alt}+\touche{↑} / \touche{Alt}+\touche{↓});
\item or click \menu{Sort by name}. The order is “natural”: \fichier{2 Flute} comes before
      \fichier{10 Violin}, as you would expect.
\end{itemize}
\menu{Remove} \ (\touche{Del}) takes the selected files out of the list; \menu{Clear all} empties
it. \textbf{No file is deleted from the disk.}

\capture[0.82]{../capture_triee-en.png}

\subsection{Checking the output file name}
Tutti suggests a name based on what the file names have in common:
\fichier{Symphony No. 5 - 01 Full score}, \fichier{Symphony No. 5 - 02 Flute}… give
\fichier{Symphony No. 5 - complete.pdf}, in the folder of the first file. You can edit it
directly, or pick another location with \menu{Browse…}.

\subsection{Merging}
Click \menu{Merge}. A bar shows the progress; the button turns into \menu{Cancel} while working.
When done, the PDF opens (option \menu{Open the PDF when done}).

\astuce{If the output file already exists, Tutti asks whether to replace it. If it is open in a
PDF reader, close it first: Windows does not allow replacing an open file.}

\section{Options}\label{sec:options}

\begin{description}[style=nextline,leftmargin=5mm,font=\sffamily\bfseries\color{bordeaux}]
\item[One bookmark per file]
The merged PDF contains one bookmark per original file (“02 Flute”, “10 Violin I”…), shown in the
side panel of your PDF reader. The bookmark is the file name without the part common to all.
\item[Start each part on a right-hand page]
For double-sided printing: when a part ends on a right-hand page, a blank page is inserted so that
the next part also starts on a right-hand page.
\item[Layout: booklet]
For a \textbf{booklet folded and stapled in the middle} (a concert programme, for example). The
booklet is made from \textbf{everything in the list}, in list order: one multi-page file, or several
one-page files (page 1, page 2…), or a mix. Blank pages are added up to a multiple of 4; they go at the
end of the booklet. Three choices:
\begin{itemize}
\item \textbf{A5 booklet ready to print}: two pages reduced to A5, side by side on each side of a
      landscape A4 sheet. Print on A4, landscape, \textbf{actual size (100\%)},
      \textbf{double-sided, flip on the short edge}.
\item \textbf{A4 booklet on A3 sheets}: the usual size for bound sheet music. Two A4 pages,
      \textbf{not reduced}, side by side on each side of a landscape A3 sheet. You need an A3
      printer (or a print shop): A3, landscape, actual size, double-sided short edge. You get an
      A4 booklet.
\item \textbf{Booklet order only}: pages keep their size and are simply put in booklet order
      (8, 1, 2, 7, 6, 3, 4, 5 for 8 pages). When printing, choose \textbf{2 pages per sheet},
      landscape, double-sided short edge.
\end{itemize}
Then fold the stack in the middle, nest the sheets and staple along the fold.
\item[Open the PDF when done]
Opens the result in your usual PDF reader. When unchecked, Tutti offers to open the PDF or its folder.
\end{description}

\section{When a file causes trouble}

A damaged or password-protected file is shown \textcolor{bordeaux}{in red} in the list, with the
reason, \textbf{before} merging. The \menu{Merge} button stays greyed out while it is there: select
it and click \menu{Remove}.

\capture[0.82]{../capture_erreur-en.png}

\begin{center}
\renewcommand{\arraystretch}{1.35}
\begin{tabularx}{\linewidth}{>{\raggedright\arraybackslash}p{0.36\linewidth} X}
\toprule
\textbf{Message} & \textbf{What to do} \\
\midrule
Unreadable or damaged PDF & Reopen the original document and export it to PDF again. \\
Password protected & Save a copy without a password from your PDF reader. \\
This PDF has no pages & The file is empty: remove it. \\
The file is probably open in another program & Close the merged PDF in your reader, then try again. \\
\bottomrule
\end{tabularx}
\end{center}

\section{Converting images and documents to PDF}\label{sec:convertir}

In the \menu{Merge and convert} tab, the list also accepts:
\begin{center}
\renewcommand{\arraystretch}{1.3}
\begin{tabularx}{\linewidth}{>{\bfseries}l X}
\toprule
Images & JPEG, PNG, TIFF (including multi-page), BMP, GIF, WebP \\
Word & \fichier{.doc}, \fichier{.docx}, \fichier{.odt}, \fichier{.rtf} \\
Excel & \fichier{.xls}, \fichier{.xlsx}, \fichier{.ods} \\
PowerPoint & \fichier{.ppt}, \fichier{.pptx}, \fichier{.odp} \\
\bottomrule
\end{tabularx}
\end{center}

Each file is converted as soon as it enters the list; it shows “(image)” or “(converted document)”
and its page count. You can then order it and merge it with PDF files like any other file.

\textbf{To convert a single file}, drop it alone into the list: Tutti suggests the same name with
the \fichier{.pdf} extension (\fichier{Letter.docx} → \fichier{Letter.pdf}), and \menu{Merge}
produces the PDF.

\begin{itemize}
\item \textbf{Images}: one image per page, on A4 in the image's orientation, without distortion.
      A JPEG photo is inserted as is: no loss of quality.
\item \textbf{Word, Excel, PowerPoint}: the layout is done by \textbf{Microsoft Office} when it is
      installed (Windows), otherwise by \textbf{LibreOffice} (free, \fichier{libreoffice.org}), required
      on Mac and Linux. It works invisibly, read-only, with macros disabled. If neither is installed,
      the file is shown in red.
\end{itemize}

\section{Touching up pages (Pages tab)}\label{sec:pages}

The \menu{Pages} tab shows \textbf{every page of a document as a thumbnail}, so you can rotate them,
remove them or change their order. It works on \textbf{one document at a time}.

\subsection{Opening the document}
Drop \textbf{one} file into the large frame, or click \menu{Open a document…}. A PDF, but also a
Word file, a PowerPoint file, an image… (converted first, as in the first tab). The thumbnails
appear one after the other; under each one, the \textbf{page number in the original document}.

\capture[0.82]{../capture_pages-en.png}

\subsection{Selecting pages}
\begin{itemize}
\item a \textbf{click} selects a page (it is highlighted);
\item \touche{Ctrl}+click adds a page to, or removes it from, the selection; \touche{Shift}+click
      selects a whole run of pages;
\item \menu{Select all} (\touche{Ctrl}+\touche{A}) selects every page;
\item a \textbf{double-click} shows the page full size, so you can read it before deciding.
\end{itemize}

\subsection{Rotate, remove, move}
\begin{description}[style=nextline,leftmargin=5mm,font=\sffamily\bfseries\color{bordeaux}]
\item[↺ Rotate left, ↻ Rotate right]
A quarter turn per click (\touche{Ctrl}+\touche{L} / \touche{Ctrl}+\touche{R}), on every selected
page. A rotated page shows the ↻ sign under its thumbnail. The page content is not touched: only its
orientation changes, with no loss of quality.
\item[Remove]
Leaves the selected pages out of the result (\touche{Del}). The original document keeps all its pages.
\item[Move]
\textbf{Drag} a thumbnail (or several selected thumbnails) to its new place.
\item[Undo all]
Goes back to the original pages, in the original order.
\end{description}
The line under the thumbnails sums up the changes: “12 pages — 1 removed — 2 rotated — order changed”.

\subsection{Saving}
\menu{Save…} asks where to write the new PDF; Tutti suggests \fichier{<name> - touched up.pdf}
next to the document. \textbf{The original document is never modified}: Tutti refuses to write over
it. If you close Tutti or open another document while changes are not saved, Tutti asks first.

\astuce{A scanned score with one page that came out sideways: open it in \menu{Pages}, click the
sideways page, \menu{Rotate}, then \menu{Save…}. Thirty seconds.}

\section{Finishing: numbers, title, size}\label{sec:finitions}

The \menu{Finishing} tab adds the final touch to one or \textbf{several} PDF files at once — all the
parts of a score, for example. Each file gives a \textbf{new} file, \fichier{<name> - finished.pdf},
next to the original (or in the folder of your choice).

\capture[0.82]{../capture_finitions-en.png}

\begin{description}[style=nextline,leftmargin=5mm,font=\sffamily\bfseries\color{bordeaux}]
\item[Page numbers]
\begin{itemize}
\item \textbf{Where}: bottom centre; or \textbf{outer edge}, i.e.\ on the right of odd pages and on
      the left of even pages — the free edge of a bound double-sided copy, where the thumb turns the
      page; at the bottom or at the top.
\item \textbf{Style}: “3”, or “3 / 12” (number and total number of pages).
\item \textbf{Start at}: the number of the first page, for a part that follows another one.
\item \textbf{No number on page 1}: the title page is not numbered, but it counts (the next one
      shows 2).
\end{itemize}
\item[Text at the top of the page]
The title of the work, the name of the part (“Smetana — The Moldau — Violin I”)… centred at the top
of \textbf{every page}, or of \textbf{page 1 only}. Leave the box empty to write nothing.
\item[Make smaller for e-mail]
Colour or grey photos and scans that are more detailed than needed are reduced to 150 dots per inch —
plenty for the screen and an office printer. Text and music notes written by software (Finale,
Sibelius, Word…) are not images: they stay perfectly sharp. Black-and-white scans are already small
and are left as they are. At the end, Tutti shows the size before and after.
\end{description}

\begin{center}
\setlength{\fboxsep}{0pt}\setlength{\fboxrule}{0.6pt}%
\fcolorbox{summit}{white}{\includegraphics[width=0.42\linewidth]{../exemple_finitions-en.png}}\\[1mm]
{\small\sffamily\color{summit} Result: title at the top, “2 / 9” at the bottom on the outer edge (even page: on the left).}
\end{center}

\astuce{Numbers and text are written in small type (10 points), 7 mm from the edge: beyond the
margin a printer cannot print, and usually above the first staff. If a score fills its page right up
to the edge, choose numbers “at the top”, or no text.}

\section{Splitting a PDF}\label{sec:decouper}

\menu{Split} tab:
\begin{enumerate}
\item drop the file(s) to split into the list: PDF files, but also Word, Excel, PowerPoint
      documents or images (converted automatically, as in the other tab);
\item choose \menu{One page per file}, or \menu{By page ranges}, typing the ranges separated by
      commas: \fichier{1-8, 9-16, 17-} (“\fichier{17-}” means “from page 17 to the end”);
\item choose where to save the files: by default, a \fichier{<name> - pages} folder is created
      next to each PDF;
\item click \menu{Split}.
\end{enumerate}
The files are named \fichier{Score - page 01.pdf}, \fichier{Score - page 02.pdf}… or
\fichier{Score - pages 01-08.pdf}: the leading zeros keep them in order in Explorer.
The original PDF is not modified. If files with the same name already exist, Tutti asks before
replacing them.

\capture[0.82]{../capture_decoupe-en.png}

\subsection*{Combining the functions}
\begin{itemize}
\item \textbf{A 6-page Word document into 6 one-page PDF files}: drop the \fichier{.docx} into the
      \menu{Split} tab and choose \menu{One page per file}.
\item \textbf{Photos and PDF files in one document}: drop them together into \menu{Merge and convert},
      put them in order, merge.
\item \textbf{An Excel sheet to PDF}: drop it alone into \menu{Merge and convert}.
\item \textbf{Pages 3 to 5 of a presentation}: drop the \fichier{.pptx} into \menu{Split}, ranges \fichier{3-5}.
\item \textbf{A complete score, clean and ready to send}: \menu{Merge} the parts → open the result in
      \menu{Pages} to straighten a sideways page and remove a blank page → \menu{Finishing}: numbers
      “3 / 12” on the outer edge, title at the top, “make smaller” → send the \fichier{- finished.pdf} file.
\item \textbf{A concert programme booklet}: merge the pages (Word, images…), then the \menu{A5 booklet}
      layout — or, to number them first, \menu{Finishing}, then \menu{Merge} the numbered file as a booklet.
\end{itemize}

\section{Updates}

At startup Tutti quietly checks for a new version. If there is one, it offers to install it: it
downloads it, installs it and restarts. Nothing has to be downloaded by hand. You can also check
yourself: \menu{Help} → \menu{Check for updates}.

\section{Reporting a problem}

\menu{Help} → \menu{Report a problem…}: describe what happened, then \menu{Send}. Without an
Internet connection, the report is saved to a file whose location Tutti shows you.

\section{Privacy}

\begin{itemize}
\item Tutti only goes online to check for updates and, with your consent, to send error reports.
\item None of your files is ever sent. In reports, file names and the user name are hidden.
\end{itemize}

\section{Keyboard shortcuts}

\begin{center}
\renewcommand{\arraystretch}{1.35}
\begin{tabular}{ll}
\toprule
\multicolumn{2}{l}{\sffamily\bfseries\color{bordeaux} Merge and convert} \\
\touche{Ctrl}+\touche{O} & Add files \\
\touche{Alt}+\touche{↑} \ / \ \touche{Alt}+\touche{↓} & Move the selected file up / down \\
\touche{Del} & Remove the selected file from the list \\
\midrule
\multicolumn{2}{l}{\sffamily\bfseries\color{bordeaux} Pages} \\
\touche{Ctrl}+\touche{L} \ / \ \touche{Ctrl}+\touche{R} & Rotate left / right \\
\touche{Del} & Remove the selected pages \\
\touche{Ctrl}+\touche{A} & Select every page \\
double-click & See the page full size \\
\midrule
\multicolumn{2}{l}{\sffamily\bfseries\color{bordeaux} Everywhere} \\
\touche{F1} & Open this guide \\
\touche{Ctrl}+\touche{Q} & Quit \\
\bottomrule
\end{tabular}
\end{center}

\vfill
{\footnotesize\color{summit}\sffamily
Tutti uses qpdf (through pikepdf), QtPdf for the thumbnails, and the Source Sans 3 and EB Garamond fonts (SIL Open Font
License). Examples in this guide use fictitious files.\par}

\end{document}
